\section{Finite driver model for direct supply to gate coupling}
The bounded change adds a capacitor $C_{FF}$ directly between $V_{IN}$
and the PFET pass gate $G$. The capacitor is in parallel with the native
gate to source coupling. It carries no DC current and does not impose
an ideal voltage at the gate. The unchanged transistor circuit must
still supply its transient charge and retain a stable connected response.

For an interpretable two node reduction, represent the existing driver
by its finite resistance $R_D$ and voltage command
\begin{equation}
v_D=K_Iv+A(s)(o-r),\qquad A(s)=\frac{A_0}{1+s\tau_A}.
\end{equation}
The positive output sensitivity is appropriate to a PFET pass device.
An output increase raises its gate and reduces delivered current.
$K_I$ is the supply sensitivity of this reduced driver. The model retains
reciprocal pass capacitors at their physical terminals and defines
\begin{align}
C_I&=C_{gs,P}+C_{FF},& C_G&=C_I+C_{gd,P},\\
C_O'&=C_{OUT}+C_{ds,P}+C_{gd,P},&
R_\Sigma&=(1/R_L+1/r_{o,P})^{-1}.
\end{align}
An ideal current load has $1/R_L=0$. Other output loading must be added
to the model when required. In particular, the actual local follower
loads both $G$ and $O$, so this two node reduction explains the capacitor
mechanism and does not replace the complete native core.

The two terminal equations are
\begin{align}
\left(1/R_D+sC_G\right)g
 -\left(A(s)/R_D+sC_{gd,P}\right)o
 &=\left(K_I/R_D+sC_I\right)v-A(s)r/R_D,\\
\left(g_{m,P}-sC_{gd,P}\right)g
 +\left(1/R_\Sigma+sC_O'\right)o
 &=\left(g_{m,P}+1/r_{o,P}+sC_{ds,P}\right)v-i_L.
\end{align}
The pass $C_{gd}$ appears in both equations. Omitting its output current
would change both the supply zero and the feedback zero.

With $F_A=1+s\tau_A$, elimination gives
\begin{align}
\frac{o}{v}&=\frac{n_0+n_1s+n_2s^2+n_3s^3}
 {d_0+d_1s+d_2s^2+d_3s^3},\\
\frac{o}{r}&=\frac{A_0g_{m,P}-A_0C_{gd,P}s}
 {d_0+d_1s+d_2s^2+d_3s^3},\\
\frac{o}{i_L}&=-\frac{1+(\tau_A+R_DC_G)s+\tau_AR_DC_Gs^2}
 {d_0+d_1s+d_2s^2+d_3s^3}.
\end{align}
The complete denominator coefficients are
\begin{align}
d_0&=1/R_\Sigma+A_0g_{m,P},\\
d_1&=C_O'+R_DC_G/R_\Sigma+R_DC_{gd,P}g_{m,P}
      +\tau_A/R_\Sigma-A_0C_{gd,P},\\
d_2&=R_D(C_GC_O'-C_{gd,P}^2)
      +\tau_A(C_O'+R_DC_G/R_\Sigma+R_DC_{gd,P}g_{m,P}),\\
d_3&=\tau_AR_D(C_GC_O'-C_{gd,P}^2).
\end{align}
Define three physical supply numerator coefficients
\begin{align}
v_0&=g_{m,P}(1-K_I)+1/r_{o,P},\\
v_1&=C_{ds,P}+K_IC_{gd,P}
      +R_D\left[g_{m,P}C_{gd,P}+C_G/r_{o,P}\right],\\
v_2&=R_D(C_GC_{ds,P}+C_IC_{gd,P}).
\end{align}
Their expanded combination with the finite feedback pole is
\begin{align}
n_0&=v_0,\qquad n_1=v_1+\tau_Av_0,\\
n_2&=v_2+\tau_Av_1,\qquad n_3=\tau_Av_2.
\end{align}
Both $C_G$ and $C_I$ contain $C_{FF}$, so the added capacitor changes
the supply numerator and the poles. It does not improve rejection by
changing only a numerator term. For this cubic model, positivity of
all denominator coefficients and $d_1d_2>d_0d_3$ are the Hurwitz
conditions. The complete transistor circuit requires its own pole
calculation because its internal node count is larger.

For every finite $C_{FF}$, the capacitor is open at DC. Therefore the
DC operating point and the exact zero frequency transfer remain
unchanged. At frequencies where $C_{FF}$ dominates the finite gate
impedance, $g$ approaches $v$. The resulting pass current is
\begin{equation}
i_P=\left[1/r_{o,P}+s(C_{ds,P}+C_{gd,P})\right](v-o).
\end{equation}
The transconductance ripple is suppressed, but finite output resistance
and drain charge still transmit current. In the limiting two node case,
\begin{equation}
\frac{o}{v}\longrightarrow
\frac{1/r_{o,P}+s(C_{ds,P}+C_{gd,P})}
 {1/R_\Sigma+s(C_{OUT}+C_{ds,P}+C_{gd,P})}.
\end{equation}
The required gate ripple for complete cancellation at fixed output is
instead
\begin{equation}
\frac{g}{v}=1+\frac{1/r_{o,P}+s(C_{ds,P}+C_{gd,P})}
 {g_{m,P}-sC_{gd,P}}.
\end{equation}
Passive common mode coupling alone does not impose this excess gate
ripple. It can also reduce reference control and increase event charge
by the same finite impedance mechanism.

The planned native cases are $C_{FF}=0,5,20,100,500$ pF. They use the
same transistor dimensions, ideal bias ports, 50 mA current load,
200 pF output capacitor and 100 pF command capacitor. Each case includes
supply and reference AC responses, the DC point and a closed pole
diagnostic. The fixed comparison frequencies are 1 kHz, 10 kHz,
100 kHz, 1 MHz, 3 MHz and 10 MHz. A pole diagnostic is accepted only if
the log is free of abort or error messages and a nonempty pole list is
returned. These tests compare a physical added capacitance rather than
an imposed gate tracking source. No novelty claim is made.
